\section*{Latihan dan Refleksi}
\addcontentsline{toc}{section}{Latihan dan Refleksi}

\subsection*{Latihan}
\begin{enumerate}[label=\arabic*.,leftmargin=*]
  \item Buktikan dengan induksi bahwa \(|w^{R}| = |w|\) untuk setiap
        string \(w\).
  \item Diberikan \(L_1 = \{a, aa\}\) dan \(L_2 = \{b, ab\}\).
        Tentukan \(L_1L_2\), \(L_1^{*}\), dan \(L_1 \cup L_2\).
  \item Tuliskan tata bahasa bebas konteks untuk bahasa
        \(\{a^i b^j \mid i \ge j \ge 0\}\).
  \item Buat derivasi leftmost dan pohon urai untuk string
        \(\code{id} + \code{id} * \code{id}\) pada tata bahasa
        aritmetika.
  \item Jelaskan mengapa bahasa \(\{a^n b^n \mid n \ge 0\}\) tidak
        dapat dikenali oleh DFA.
  \item Berikan contoh bahasa sensitif konteks yang bukan bebas
        konteks, lalu jelaskan mengapa satu stack tidak cukup.
  \item Mengapa analisis leksikal cukup memakai otomata hingga,
        sedangkan analisis sintaksis memerlukan otomata pushdown?
  \item Jelaskan hubungan antara pohon urai dan AST. Mengapa AST lebih
        ringkas?
\end{enumerate}

\subsection*{Refleksi}
\begin{reflectionbox}
\begin{enumerate}[label=\arabic*.,leftmargin=*]
  \item Konsep formal mana yang paling sulit Anda pahami? Mengapa?
  \item Apakah Anda dapat menjelaskan perbedaan bahasa regular dan
        bebas konteks tanpa membaca catatan?
  \item Bagaimana notasi formal membantu Anda menulis spesifikasi
        bahasa yang lebih tepat?
\end{enumerate}
\end{reflectionbox}
